\section{Related Work}
\label{sec-related}

\stitle{Natural language interfaces to data}.
Natural language interfaces to relational databases map a parsed
question onto SQL, and NaLIR made the mapping interactive by asking the
user to resolve ambiguities~\cite{li2014nalir}.
For RDF, gAnswer casts question answering as subgraph matching between
a semantic query graph and the knowledge graph~\cite{zou2014natural},
and template decomposition answers complex questions by composing
binary templates mined from the graph~\cite{zheng2018question}.
GraphLingo explores domain knowledge graphs in natural language by
synchronizing them with an LLM~\cite{le2024graphlingo}.
Benchmarks for question answering over RDF include LC-QuAD
2.0~\cite{dubey2019lcquad}, QALD-9-plus~\cite{qald9plus},
QALD-10~\cite{qald10}, and the Text2SPARQL
challenges~\cite{text2sparql2025,text2sparql}.

\stitle{LLM-based text-to-SQL}.
With LLMs, text-to-SQL has moved to prompting and fine-tuning.
DAIL-SQL studies how question representation and example selection
shape LLM prompts~\cite{gao2024dailsql}, PURPLE retrieves demonstrations
that match the logical operators a question needs~\cite{ren2024purple},
CodeS pre-trains open models for SQL generation~\cite{li2024codes}, and
NL2SQL360 evaluates these systems along many dimensions~\cite{li2024dawn}.
These methods emit a whole query and improve it by better context or
better models; none types the query against the schema before it runs.

\stitle{LLM-based text-to-SPARQL}.
SPARQL-LLM, the strongest entry of the 2026 Text2SPARQL challenge,
generates SPARQL from natural-language questions in real
time~\cite{smeros2026sparqlllm}.
ARUQULA explores the graph with tools in a ReAct loop before writing a
query~\cite{brei2025aruqula}.
UniQGen extracts a table of constraints from the question and the graph
and refines a query with Chase and Backchase against reference answers
from an LLM, rendering it to SPARQL or Cypher~\cite{wang2026uniqgen}.
\hermes differs from these in its unit of generation: an operator of
the SPARQL algebra rather than a query string or a constraint table,
checked against the schema before execution and observed after it.

\stitle{Grounding and selection}.
KB-BINDER drafts a logical form with an LLM and binds its terms to the
knowledge base by lexical matching~\cite{li2023kbbinder}, and Pangu lets
a symbolic agent propose only valid plans while the language model scores
them~\cite{gu2023pangu}.
CHASE-SQL generates SQL candidates along several reasoning paths and
selects one with a trained pairwise selector~\cite{pourreza2025chasesql},
and ThriftLLM chooses which LLMs to query for classification under a cost
budget from their estimated success rates~\cite{huang2025thriftllm}.
The budgeted planner of Section~\ref{sec-framework} combines the two
lines: candidates are bound to the ontology and type-checked, then
chosen by an estimate of correctness net of an estimated execution
price, under a budget of model calls and executions.

\stitle{Query optimization and SPARQL semantics}.
Cost-based optimizers choose among algebraically equivalent plans by
estimating the effect of each operator~\cite{selinger1979access}, and
cardinality estimates dominate the quality of the chosen
plan~\cite{leis2015good}.
Learned components now steer existing optimizers~\cite{marcus2021bao}
and estimate the cardinality of knowledge-graph
queries~\cite{davitkova2022lmkg,schwabe2024gnce}; \hermes prices
candidate queries with classical estimates and could use either.
Chase and Backchase reformulates queries under constraints and
minimizes them~\cite{deutsch1999chase}.
Pérez et al.\ gave SPARQL a compositional semantics and characterized
the complexity of its evaluation~\cite{perez2009semantics}.
\hermes borrows the optimizer's step, one operator with an observed
effect, and the algebra's semantics, and uses them to generate a query
rather than to rewrite one.
